\newcommand{\DoNotLoadEpstopdf}{}
\pdfinfoomitdate=1
\pdftrailerid{}
\pdfsuppressptexinfo=-1
\documentclass[11pt,letterpaper]{article}
\usepackage[T1]{fontenc}
\usepackage{lmodern}
\usepackage{amsmath,amssymb,amsthm,mathtools}
\usepackage[margin=1in]{geometry}
\usepackage{booktabs,array,microtype}
\usepackage[hidelinks]{hyperref}
\hypersetup{pdftitle={Report 168 Sparse binary matrices with line sums at most two},pdfauthor={},pdfsubject={Exact generating functions all fixed order asymptotics and safe inversion of A197458}}
\setlength{\parskip}{4pt}
\setlength{\parindent}{1.2em}
\setlength{\emergencystretch}{3em}
\setcounter{tocdepth}{1}
\numberwithin{equation}{section}
\newtheorem{theorem}{Theorem}[section]
\newtheorem{lemma}[theorem]{Lemma}
\newtheorem{proposition}[theorem]{Proposition}
\newtheorem{corollary}[theorem]{Corollary}
\theoremstyle{remark}\newtheorem{remark}[theorem]{Remark}
\newcommand{\dd}{\,\mathrm d}
\newcommand{\CC}{\mathbb C}
\newcommand{\QQ}{\mathbb Q}
\newcommand{\NN}{\mathbb N}
\newcommand{\Rea}{\operatorname{Re}}
\newcommand{\Imma}{\operatorname{Im}}
\newcommand{\ceil}[1]{\left\lceil #1\right\rceil}
\title{Sparse binary matrices with line sums at most two\\\large Exact counts and all fixed order asymptotics}
\author{Report 168}
\date{3 October 2026}
\begin{document}
\maketitle
\begin{abstract}
Let $a_n$ count $n\times n$ binary matrices with at most two ones in
every row and column, with $a_0=1$. This is OEIS A197458. We prove
\[
 a_n\sim \frac{(n!)^2}{2^{5/4}\pi n^{3/4}}e^{2\sqrt{2n}},
\]
and an expansion to every fixed order in $n^{-1/2}$, with four explicit
corrections and a finite exact Gaussian-moment rule for all subsequent
coefficients. The proof starts from an elementary labeled-component
decomposition and a modified-Bessel diagonal. A sectorial Bessel
remainder, a whole-circle bound, explicit off-circle connector
estimates, and an integrable Taylor remainder give the claimed sharp
fixed-order error.

We also derive a logarithmic expansion, an explicit Lambert-$W$
real-inverse formula through its bounded correction, and a rigorous
two-ceiling enclosure for the integer threshold. This enclosure
retains the distinction between a smooth comparison inverse and the
actual discrete crossing. The companion supplies bounded exact counts,
independent finite checks, and optional symbolic coefficient checks.

General bounded-degree enumeration and configuration-model methods are
classical. This is an explicit specialization and refinement; no
historical-priority assertion is made. The results are fixed-order
asymptotics, not a complete beyond-all-orders description. Remainder
constants and eventual thresholds are existential, not numerically
certified.
\end{abstract}
\newpage
\tableofcontents
\medskip
\noindent\textbf{Reading guide.}
Section~\ref{sec:results} states the asymptotic results. Sections
\ref{sec:exact}--\ref{sec:gaussian} prove the exact formulas and the
coefficient expansion. Section~\ref{sec:inverse} proves real and integer
inversion. Sections~\ref{sec:checks}--\ref{sec:literature} record
reproducibility and the precisely bounded source coverage. All analytic
arguments used for the main theorem are included.
\newpage

\section{Results and conventions}\label{sec:results}
Rows and columns are individually labeled. A line means either a row
or a column. Write $b_n=a_n/(n!)^2$, and use $[z^k]$ for coefficient
extraction. Every asymptotic statement below is as its argument tends
to positive infinity. An order index such as $J$ is fixed before the
limit; implied constants may depend on it.

\begin{theorem}\label{thm:main}
For each integer $J\geq0$ there are constants $c_j\in\QQ(\sqrt2)$ such
that
\begin{equation}\label{eq:main}
 a_n=\frac{(n!)^2e^{2\sqrt{2n}}}{2^{5/4}\pi n^{3/4}}
 \left(\sum_{j=0}^J c_jn^{-j/2}
              +O_J\bigl(n^{-(J+1)/2}\bigr)\right).
\end{equation}
The first five coefficients are
\begin{align}\label{eq:cs}
 c_0&=1,&c_1&=-\frac{17\sqrt2}{96},&c_2&=\frac{6409}{9216},\notag\\
 c_3&=-\frac{18114133\sqrt2}{13271040},&
 c_4&=\frac{13052721749}{2548039680}.
\end{align}
The finite rule in Section~\ref{sec:gaussian} specifies every $c_j$.
\end{theorem}
The theorem asserts a Poincar\'e expansion at every fixed order. It
neither asserts convergence of the formal coefficient series nor
identifies every exponentially small contribution. In particular,
letting $J$ grow with $n$ is not justified by the theorem.

\begin{corollary}\label{cor:log}
The logarithmic expansion begins
\begin{align}\label{eq:log}
 \log a_n={}&2n(\log n-1)+2\sqrt{2n}
       +\tfrac14\log n-\tfrac14\log2\notag\\
 &-\frac{17\sqrt2}{96\sqrt n}+\frac{319}{384n}
       +O(n^{-3/2}).
\end{align}
There is a logarithmic expansion to every fixed half-integer order.
\end{corollary}
For the latter assertion, formally expand
$\log(\sum c_jn^{-j/2})$ to the desired finite order and add twice
Stirling's logarithmic series. The constant and first correction check
is
\[
 c_2-\tfrac12c_1^2=\frac{85}{128},\qquad
 \frac{85}{128}+\frac16=\frac{319}{384}.
\]
Positivity of the leading normalized quantity for sufficiently large
$n$ justifies taking its logarithm with the stated remainder.

\section{Exact component decomposition and diagonal}\label{sec:exact}
A matrix is a simple bipartite graph with specified shores of size $n$.
Maximum degree two means that each connected component is an isolated
vertex, a path, or an even cycle. Let $A_{m,n}$ count the rectangular
class and define its bivariate exponential generating function
\[
 F(x,y)=\sum_{m,n\geq0}A_{m,n}\frac{x^m y^n}{m!n!}.
\]
For $t=xy$, the connected component contributions are as follows:
\begin{center}
\begin{tabular}{ll}
\toprule
Component & Contribution\\
\midrule
Isolated vertices & $x+y$\\
Balanced paths on $k+k$ vertices, $k\geq1$ & $t^k$\\
Unbalanced paths on $(k+1)+k$ vertices, $k\geq1$
 & $(x+y)t^k/2$\\
Even cycles on $k+k$ vertices, $k\geq2$ & $t^k/(2k)$\\
\bottomrule
\end{tabular}
\end{center}
A balanced path has $(k!)^2$ realizations on specified labels: choose
its ordering starting from its unique endpoint in the first shore.
An unbalanced path has a reversal symmetry of order two. A cycle has
$2k$ equivalent oriented starting descriptions. The lower cutoff
$k=2$ for cycles excludes the parallel-edge component forbidden in a
binary matrix. The exponential formula for sets of labeled connected
components now gives
\begin{equation}\label{eq:egf}
 F(x,y)=(1-t)^{-1/2}
 \exp\left\{\frac{t}{1-t}-\frac t2
             +(x+y)\frac{2-t}{2(1-t)}\right\}.
\end{equation}
This identity also follows directly by multiplying the exponential
series for the four kinds of components, so no asymptotic input is
being invoked.

Put
\begin{equation}\label{eq:Pq}
 P(t)=(1-t)^{-1/2}e^{t/(1-t)-t/2},\quad
 h(t)=\frac{2-t}{2(1-t)},\quad
 q(t)=2\sqrt t\,h(t).
\end{equation}
The diagonal requires equal powers from $e^{xh(t)}$ and $e^{yh(t)}$.
Consequently
\begin{equation}\label{eq:bessel}
 G(t):=\sum_{n\geq0}b_nt^n
   =P(t)\sum_{k\geq0}\frac{t^kh(t)^{2k}}{(k!)^2}
   =P(t)I_0\left(\frac{\sqrt t(2-t)}{1-t}\right),
\end{equation}
where
\[
 I_0(z)=\sum_{k\geq0}\frac{(z^2/4)^k}{(k!)^2}.
\]
The square root is harmless because $I_0$ is even. In particular,
$q(t)^2=t(2-t)^2/(1-t)^2$ is rational, and $G$ is analytic in $|t|<1$
with nonnegative coefficients.

\section{Differential equation and exact recurrence}\label{sec:ode}
Let
\[
 R(t)=\frac{t(2-t)^2}{(1-t)^2},\quad
 L(t)=\frac{P'(t)}{P(t)}=\frac1{2(1-t)}+\frac1{(1-t)^2}-\frac12.
\]
Termwise differentiation shows that $f(R)=I_0(\sqrt R)$ solves
$Rf''+f'-f/4=0$. Under the pullback $R=R(t)$ define
\[
 A=\frac{R'}R-\frac{R''}{R'},\qquad B=\frac{R'^2}{4R}.
\]
The pulled-back solution satisfies $Y''+AY'-BY=0$. Conjugating by
$P$, then clearing denominators, yields
\begin{equation}\label{eq:ode}
 p_2(t)G''+p_1(t)G'+p_0(t)G=0,
\end{equation}
with
\begin{align}\label{eq:poly}
 p_0(t)&=t^6-2t^5-3t^4+20t^3-36t^2+32,\notag\\
 p_1(t)&=4t^6-12t^5+8t^4+28t^3-92t^2+96t-16,\\
 p_2(t)&=4t^6-20t^5+44t^4-60t^3+48t^2-16t.\notag
\end{align}
For example, the unnormalized coefficients before clearing
 denominators are $A-2L$ and $L^2-L'-AL-B$; this gives a short rational
identity check of every displayed polynomial.

Write $p_{i,k}=[t^k]p_i(t)$ and interpret negative-index $b$'s as zero.
The coefficient equation is
\begin{align}\label{eq:recurrence}
0={}&\sum_{k=0}^6p_{2,k}(n-k+2)(n-k+1)b_{n-k+2}\notag\\
 &+\sum_{k=0}^6p_{1,k}(n-k+1)b_{n-k+1}
   +\sum_{k=0}^6p_{0,k}b_{n-k}.
\end{align}
The coefficient of $b_{n+1}$ is $-16(n+1)^2$. Thus $b_0=1$ uniquely
determines the analytic power series. For a recurrence involving only
integers, multiply the other terms by $(n!)^2$ before solving for
$a_{n+1}$: since
$16(n+1)^2b_{n+1}(n!)^2=16a_{n+1}$, every remaining term is an integer
multiple of $a_m(n!/m!)^2$ with $0\leq m\leq n$. Division by $16$ is
exact, as follows either from the generating function or the
combinatorial interpretation. The companion checks this division at
every step rather than silently truncating it.

\section{A sectorial Bessel expansion with remainder}\label{sec:bessel}
The following elementary proof specifies the uniformity needed later.
For all complex $q$,
\begin{equation}\label{eq:bessel-integral}
 I_0(q)=\frac1\pi\int_0^\pi e^{q\cos v}\dd v
       =\frac{e^q}\pi\int_0^2
          e^{-qw^2/2}(1-w^2/4)^{-1/2}\dd w.
\end{equation}
The first identity follows by expanding the exponential and
integrating powers of cosine; the second uses the real substitution
$w=2\sin(v/2)$.

\begin{lemma}\label{lem:bessel}
For any fixed $M\geq1$ and $\delta>0$, uniformly as $|q|\to\infty$ in
$|\arg q|\leq\pi/2-\delta$,
\begin{equation}\label{eq:bessel-asym}
 I_0(q)=\frac{e^q}{\sqrt{2\pi q}}
  \left(\sum_{k=0}^{M-1}\gamma_kq^{-k}+O_{M,\delta}(|q|^{-M})\right),
 \qquad
 \gamma_k=\frac{\bigl(1\cdot3\cdots(2k-1)\bigr)^2}{k!8^k},
\end{equation}
where $\gamma_0=1$ and the square root is positive on the positive axis.
\end{lemma}
\begin{proof}
On $0\leq w\leq1$ expand
\[
 (1-w^2/4)^{-1/2}
   =\sum_{k=0}^{M-1}\frac{\binom{2k}{k}}{16^k}w^{2k}
       +O_M(w^{2M}).
\]
Since $\Rea q\geq |q|\sin\delta$, the integral of the remainder is
$O_{M,\delta}(|q|^{-M-1/2})$. The interval $1\leq w\leq2$ contributes
an exponentially small quantity: the endpoint singularity is
integrable and $|e^{-qw^2/2}|\leq e^{-|q|\sin\delta/2}$ there.
Extending each retained monomial integral from $[0,1]$ to
$[0,\infty)$ has the same type of exponentially small error. Finally,
\[
 \int_0^\infty e^{-qw^2/2}w^{2k}\dd w
   =2^{k-1/2}\Gamma(k+\tfrac12)q^{-k-1/2}.
\]
This is a real substitution for positive $q$, and then holds throughout
$\Rea q>0$ by holomorphic continuation; local domination makes the
integral holomorphic there. Combining factors proves the formula.
No complex substitution of the original real integration path is used.
\end{proof}

Set $s=1-t$ and
\begin{equation}\label{eq:local-p}
 p(s)=\sqrt{1-s}(1+s),\qquad q(t)=p(s)/s,
\end{equation}
where $p$ is analytic and positive at zero. Direct substitution into
\eqref{eq:bessel} gives, in a sufficiently narrow fixed sector about
positive real $s$,
\begin{equation}\label{eq:localG}
 G(1-s)=\frac{e^{-1}}{\sqrt{2\pi}}e^{2/s}
 \left(\sum_{j=0}^{M-1}h_js^j+O_M(|s|^M)\right).
\end{equation}
The formal series $H(s)=\sum h_js^j$ is specified by
\begin{equation}\label{eq:H}
 H(s)=p(s)^{-1/2}
   \exp\left(\frac{p(s)-1}{s}-\frac12+\frac s2\right)
   \sum_{k\geq0}\gamma_k\left(\frac{s}{p(s)}\right)^k.
\end{equation}
Only finite truncations are used. The first factors are analytic at
zero: the apparent singularity in $(p(s)-1)/s$ is removable. Their
product with a finite Bessel expansion gives the stated remainder, but
no convergence of the infinite formal $H$ is asserted. In particular,
\begin{align}\label{eq:hexp}
 H(s)&=1-\frac s4+\frac{7s^2}{32}-\frac{47s^3}{384}
                      +\frac{1361s^4}{6144}+O(s^5),\notag\\
 \frac{p(s)-1}{s}-\frac12+\frac s2
     &=-\frac s8-\frac{3s^2}{16}-\frac{13s^3}{128}
                      -\frac{17s^4}{256}+O(s^5).
\end{align}
Here and subsequently a finite-order $O$ attached to a formal series
means the corresponding finite truncation, justified analytically by
\eqref{eq:localG} when it is used in a contour integral.

\section{Whole-circle control and contour deformation}\label{sec:arcs}
Choose
\[
 s_0=\sqrt{2/n},\qquad r=1-s_0,\qquad
 t=re^{i\theta},\quad -\pi\leq\theta\leq\pi.
\]
We use only sufficiently large $n$, so $0<r<1$. Because
$q(t)=\sqrt t\,(1+(1-t)^{-1})$,
\[
 |q(t)|\leq\sqrt r\,(1+(1-r)^{-1})=q(r).
\]
The nonnegative coefficients of $I_0$ give
$|I_0(q(t))|\leq I_0(|q(t)|)\leq I_0(q(r))$. Also
$|1-t|\geq1-r$. Define the exact loss
\begin{equation}\label{eq:D}
 D(r,\theta)=\frac r{1-r}-\Rea\frac t{1-t}
  =\frac{r(1+r)(1-\cos\theta)}
        {(1-r)|1-re^{i\theta}|^2}.
\end{equation}
The remaining factor $e^{-t/2}$ costs at most $e$, and hence
\begin{equation}\label{eq:global}
 |G(re^{i\theta})|\leq e\,G(r)e^{-D(r,\theta)}.
\end{equation}
Uniformly on the whole circle,
\begin{equation}\label{eq:Dlower}
 D(r,\theta)\geq c\min\{\theta^2/s_0^3,1/s_0\}
\end{equation}
for an absolute $c>0$. Indeed $1-\cos\theta$ is bounded above and below
by positive multiples of $\theta^2$ on $[-\pi,\pi]$, and
$|1-re^{i\theta}|^2=s_0^2+2r(1-\cos\theta)$.

Take $\theta_0=s_0^{3/2}\log n$. By \eqref{eq:global} and
\eqref{eq:Dlower}, the coefficient integral outside
$|\theta|\leq\theta_0$ is bounded by
\begin{equation}\label{eq:minor}
 O\left(r^{-n}G(r)e^{-c\log^2n}\right).
\end{equation}
The positive-axis instance of \eqref{eq:localG} shows
$r^{-n}G(r)=O(e^{2\sqrt{2n}})$. Thus \eqref{eq:minor} is smaller than
$e^{2\sqrt{2n}}n^{-3/4}$ by every fixed power of $n$. This estimates
all other Bessel sectors without using the local Bessel expansion
there.

On the retained arc, $s=1-t$ satisfies $|s|\asymp s_0$ and
$\arg s=O(n^{-1/4}\log n)$, so \eqref{eq:localG} is uniform.
To see explicitly that its remainder retains the Gaussian width,
observe the exact circle identity
\begin{equation}\label{eq:Re2s}
 \Rea(2/s)=2/s_0-2D(r,\theta).
\end{equation}
An $O(s_0^M)$ amplitude remainder contributes at most a constant times
\[
 s_0^Mr^{-n}e^{2/s_0}
       \int_{-\theta_0}^{\theta_0}e^{-2D(r,\theta)}\dd\theta
  =O\left(s_0^{M+3/2}r^{-n}e^{2/s_0}\right).
\]
The equality follows from \eqref{eq:Dlower} and the substitution
$\theta=s_0^{3/2}v$. Relative to the saddle scale this is $O(s_0^M)$,
without any extra power or logarithm of $n$.

For a particularly simple coefficient calculator we deform the local
arc into the vertical segment
\[
 s=s_0+iv,\qquad |v|\leq v_n:=r\sin\theta_0.
\]
At each end, a horizontal connector joins this segment to the original
arc, whose real coordinate is $1-r\cos\theta_0$. Connector length is
$O(\theta_0^2)$. The region is analytic and avoids $s=0$, and stays
inside $|t|<1$ for large $n$.
Define the real phase
\[
 E(s)=\Rea\{2/s-(n+1)\log(1-s)\}.
\]
At a vertical endpoint the exact difference is
\begin{equation}\label{eq:endpoint}
 E(s_0+iv_n)-E(s_0)
  =-\frac{2v_n^2}{s_0(s_0^2+v_n^2)}
    -\frac{n+1}{2}\log(1+v_n^2/r^2)
  \leq-c\log^2n.
\end{equation}
On either connector,
\[
 \left|\frac{\partial E}{\partial\Rea s}\right|
 \leq\frac2{|s|^2}+\frac{n+1}{|1-s|}=O(n).
\]
Therefore its phase varies by only
$O(n\theta_0^2)=O(n^{-1/2}\log^2n)=o(1)$. The finite amplitude from
\eqref{eq:localG} is bounded there. The connector integrals are again
smaller than the saddle term by every power. This separate estimate
is essential: the circle-only bound \eqref{eq:global} was not applied
to an off-circle path.

\section{Gaussian coefficient rule and sharp remainder}\label{sec:gaussian}
Let $\varepsilon=n^{-1/4}$ and put
\begin{equation}\label{eq:scale}
 s=\sqrt2\,\varepsilon^2+i2^{1/4}\varepsilon^3u.
\end{equation}
The vertical endpoints satisfy $|u|\sim\sqrt2\log n$. Define
\begin{equation}\label{eq:Psi}
 \Psi(\varepsilon,u)=\frac2s-(\varepsilon^{-4}+1)\log(1-s)
                  -2\sqrt2\,\varepsilon^{-2}-1+u^2.
\end{equation}
Expanding the rational term and logarithm shows that the negative
powers cancel, and the constant term is zero. Thus
$\Psi\in\varepsilon\CC[u][[\varepsilon]]$. Notice the $n+1$ in this
expression, which comes from the coefficient kernel $(1-s)^{-n-1}$.

Define the Gaussian functional $\mathcal M$ on polynomials by
\begin{equation}\label{eq:moment}
 \mathcal M(u^{2k})=\frac{(2k)!}{4^kk!},\qquad
 \mathcal M(u^{2k+1})=0.
\end{equation}
Equivalently, $\mathcal M(Q)=\pi^{-1/2}\int_{\mathbb R}e^{-u^2}Q(u)\dd u$.
The promised all-fixed-order rule is
\begin{equation}\label{eq:calculator}
 \boxed{\quad c_j=\mathcal M\left(
 [\varepsilon^{2j}]\,H(s)e^{\Psi(\varepsilon,u)}\right).\quad}
\end{equation}
For a requested $j$, retain $H$ through degree $j$, expand $\Psi$
through degree $2j$, and perform finite polynomial multiplication.
Although $2^{1/4}$ and $i$ appear in the intermediate scaling, only
even powers survive the Gaussian moments, giving $c_j\in\QQ(\sqrt2)$.
One can also see this from the parity of the powers of $u$ in the
expansion. All operations are exact algebraic operations.

For an explicit implementation, if $\Psi=\sum_{m\geq1}\psi_m(u)
\varepsilon^m$ and $e^\Psi=\sum E_m(u)\varepsilon^m$, use
\begin{equation}\label{eq:exp-rec}
 E_0=1,\qquad E_m=\frac1m\sum_{k=1}^m k\psi_k E_{m-k}.
\end{equation}
It follows by differentiating the formal exponential. Equations
\eqref{eq:H}, \eqref{eq:Psi}, \eqref{eq:moment}, and \eqref{eq:exp-rec}
are a finite calculator at every requested fixed order, independent
of any numerical fitting.

\begin{proof}[Proof of Theorem~\ref{thm:main}]
The discarded circle arcs and horizontal connectors are already
controlled. On the vertical segment substitute \eqref{eq:localG} and
\eqref{eq:scale} into Cauchy's formula. The leading constant is
\[
 \frac{e^{-1}}{\sqrt{2\pi}}\,
 \frac{2^{1/4}}{2\pi}\,e\,\sqrt\pi
       =\frac1{2^{5/4}\pi}.
\]
The differential supplies $\varepsilon^3=n^{-3/4}$, while the
remaining exponential is $e^{2\sqrt{2n}}e^{-u^2}e^\Psi$.
In particular the $e^{-1}$ in the local Bessel expression cancels
the $e$ from coefficient extraction; dropping that cancellation
would change the leading normalization.

It remains to justify finite expansion and its sharp integrated
remainder. Choose the Bessel expansion with $M=J+1$. Its error is
$O(|s|^{J+1})=O(\varepsilon^{2J+2})$, and integrates to that relative
order by \eqref{eq:Re2s}, or by the analogous vertical Gaussian bound
from \eqref{eq:endpoint}. Use the resulting finite polynomial
$H_J(s)=\sum_{j=0}^Jh_js^j$ in the analytic factor
$H_J(s)e^{\Psi(\varepsilon,u)}$.

After cancellation in \eqref{eq:Psi}, this factor is analytic in
$\varepsilon$ near zero whenever $|\varepsilon u|$ is sufficiently
small. For each fixed derivative order, the derivatives of its
rational and logarithmic factors are bounded by fixed polynomials
in $|u|$, uniformly for $|u|\leq C\log n$ and every real intermediate
$0\leq\varepsilon'\leq\varepsilon$. The canceled phase obeys
\[
 |\Psi(\varepsilon',u)|
       \leq C\varepsilon'(1+|u|+|u|^3).
\]
Indeed this follows from the convergent geometric expansion of
$(1+i2^{-1/4}\varepsilon' u)^{-1}$ and the logarithm expansion in
$s$, after removing the finitely many Laurent and constant terms.
Since $\varepsilon\log n\to0$, the cubic term is at most $u^2/4$
for large $n$, while the other terms are bounded by $u^2/4+O(1)$.
Taylor's integral remainder through degree $2J+1$, multiplied by
$e^{-u^2}$, is therefore bounded by
\begin{equation}\label{eq:gaussian-rem}
 C_J\varepsilon^{2J+2}(1+|u|)^{D_J}e^{-u^2/2}
\end{equation}
for a fixed $D_J$. Its integral is $O_J(\varepsilon^{2J+2})$,
with no logarithmic loss.

The identity $s(-\varepsilon,-u)=s(\varepsilon,u)$, and the identical
symmetry for $\Psi$, makes every odd $\varepsilon$ coefficient an odd
polynomial in $u$. Its integral over the symmetric segment vanishes.
The Gaussian tails of each retained polynomial outside the segment
are superpolynomially small. Thus the even coefficients are precisely
\eqref{eq:calculator}, proving \eqref{eq:main} and its remainder.
The finite calculation through degree eight yields \eqref{eq:cs}.
\end{proof}

\subsection*{An independent algebraic check of the corrections}
The ODE recurrence gives a useful second route to the first four
coefficients, without proving existence of the expansion. Insert
\[
 b_n=C e^{2\sqrt{2n}}n^{-3/4}\sum_{j\geq0}c_jn^{-j/2}
\]
into \eqref{eq:recurrence}, put $z=n^{-1/2}$, and normalize each shift
$n\mapsto n+d$ by the unshifted leading factor. Its contribution is
\[
 \exp\left(\frac{2\sqrt2}{z}
                  (\sqrt{1+dz^2}-1)\right)
 \sum_{j\geq0}c_jz^j(1+dz^2)^{-3/4-j/2}.
\]
Expand the recurrence after multiplying by $z^4$. Successive nonzero
equations, with previous coefficients substituted, are
\begin{align*}
 8\sqrt2c_1+17/6&=0,\\
 \sqrt2(9216c_2-6409)/576&=0,\\
 24\sqrt2c_3+18114133/276480&=0,\\
 \sqrt2(2548039680c_4-13052721749)/79626240&=0.
\end{align*}
They independently recover \eqref{eq:cs}. This is a finite symbolic
identity check, not a replacement for the all-arcs and remainder proof.

\section{Real inverse and discrete threshold}\label{sec:inverse}
For a fixed $J\geq0$ define the explicit smooth comparison
\begin{equation}\label{eq:FJ}
 F_J(x)=\frac{\Gamma(x+1)^2e^{2\sqrt{2x}}}
                   {2^{5/4}\pi x^{3/4}}
                \sum_{j=0}^Jc_jx^{-j/2}.
\end{equation}
The correction polynomial tends to one, so $F_J$ is positive for
all sufficiently large $x$. Differentiated Stirling expansions give
\begin{equation}\label{eq:derivative}
 (\log F_J)'(x)=2\log x+O_J(x^{-1/2}),
\end{equation}
and therefore it is eventually strictly increasing and unbounded.
Let $x_J(y)$ denote its inverse on any such eventual interval.
The interval's finite initial choice has no effect for large $y$.

\begin{theorem}\label{thm:smooth-inverse}
For $L=\log y$, let $W$ be the positive real branch of Lambert's
function and put
\begin{equation}\label{eq:N}
 N=\frac{L}{2W(L/(2e))},\qquad q=\log N.
\end{equation}
For every fixed $J$, the inverse has the expansion
\begin{equation}\label{eq:inverse}
 x_J(y)=N-\frac{\sqrt{2N}}q-\frac18+
             \frac{\log2}{8q}+\frac1{q^2}-\frac1{q^3}
             +O_J\left(\frac1{\sqrt N\,q}\right).
\end{equation}
Moreover,
\begin{equation}\label{eq:sequence-inverse}
 x_J(a_n)=n+O_J\left(\frac{n^{-(J+1)/2}}{\log n}\right).
\end{equation}
In particular the right side of \eqref{eq:inverse}, evaluated with
$y=a_n$, equals $n+O(1/(\sqrt{N}\log N))$.
\end{theorem}
\begin{proof}
The definition of $W$ gives exactly $2N(\log N-1)=L$.
For $f(x)=2x(\log x-1)$ we have $f'(N)=2q$ and $f''(N)=2/N$.
Stirling's formula in \eqref{eq:FJ} gives
\[
 \log F_J(x)=f(x)+2\sqrt{2x}
                  +\tfrac14\log x-\tfrac14\log2+O_J(x^{-1/2}).
\]
Set $\delta_0=-\sqrt{2N}/q$. Taylor expansion at $N$ shows that
$\log F_J(N+\delta_0)-L$ equals
\begin{equation}\label{eq:inverse-residual}
 \frac{q-\log2}{4}-\frac2q+\frac2{q^2}+O_J(N^{-1/2}).
\end{equation}
The $\sqrt N$ terms cancel by the choice of $\delta_0$; the quadratic
term of $f$ contributes $2/q^2$, and the linear term of the square
root contributes $-2/q$.
Dividing the displayed part of \eqref{eq:inverse-residual} by $-2q$
gives
\[
 \delta_1=-\frac18+\frac{\log2}{8q}+\frac1{q^2}-\frac1{q^3}.
\]
This is bounded. The cross term of $f$ involving $\delta_0\delta_1$
is $O(1/(\sqrt N q))$; its cubic remainder is
$O(1/(\sqrt Nq^3))$. The square-root term involving $\delta_1$ is
$O(N^{-1/2})$, and its quadratic remainder is
$O(1/(\sqrt Nq^2))$. The logarithmic term changes by
$O(1/(\sqrt Nq))$. These estimates prove that
\[
 \log F_J(N+\delta_0+\delta_1)-L=O_J(N^{-1/2}).
\]
The derivative in \eqref{eq:derivative} is at least $q$ near this point
for large $N$. For completeness, the residual and this derivative
bound place a root inside a fixed sufficiently large multiple of
$1/(\sqrt Nq)$ of the point by the intermediate value theorem.
Monotonicity identifies that root as $x_J(y)$, proving
\eqref{eq:inverse}.

By Theorem~\ref{thm:main},
$\log a_n-\log F_J(n)=O_J(n^{-(J+1)/2})$.
Bracketing near $n$ with derivative at least $\log n$, or applying
the mean value theorem after that bracketing, proves
\eqref{eq:sequence-inverse}. The final statement follows by combining
this estimate with \eqref{eq:inverse}; $N\sim n$.
\end{proof}

\subsection*{Constructing further fixed inverse orders}
All higher fixed inverse accuracy is constructive. One may use the
known finite expansion of $\log F_J$ and Newton reversion, with the
error divided by a derivative bounded below by $\log x$. Here is an
explicit procedure with a quantitative stopping criterion.
For any fixed $B>0$, Stirling's series and the finite logarithmic
expansion of the correction polynomial give an explicit elementary
function $T(x)$ with
\[
 T(x)=\log F_J(x)+O_{J,B}(x^{-B}),\quad
 T'(x)=2\log x+O_J(x^{-1/2}),\quad T''(x)=O_J(1/x).
\]
Retain sufficiently many half-powers (and Stirling terms) to ensure
that remainder. The derivative statements follow by differentiating
the finite expression, rather than differentiating an unspecified
$O$ term. The corresponding remainder for the original smooth
function follows from the usual Stirling remainder and the convergent
logarithm of the finite correction polynomial near one.

Starting at $X_0=N$, iterate the entirely explicit rule
\begin{equation}\label{eq:newton}
 X_{k+1}=X_k-\frac{T(X_k)-L}{T'(X_k)}.
\end{equation}
The root of $T(x)=L$ is within $O(\sqrt N/\log N)$ of $N$ and within
$O_{J,B}(N^{-B}/\log N)$ of $x_J(y)$. On that neighborhood the Newton
error satisfies
\[
 |e_{k+1}|\leq\frac{C_J}{N\log N}|e_k|^2.
\]
This follows from Taylor's theorem, the bound on $T''$, and the lower
bound on $T'$. For every fixed desired power of $N^{-1}$, a finite
number of iterations suffices. Choosing $B$ larger than that desired
power yields a finite expression to that accuracy. Expanding these
finite expressions gives the customary formal inverse series.
For an inverse valid at sequence values, increasing $J$ in
\eqref{eq:sequence-inverse} supplies any requested fixed half-power
accuracy divided by $\log n$. This procedure does not determine
unspecified effective error constants or threshold locations.

\begin{theorem}\label{thm:integer}
Define the actual integer crossing by
\[
 I(y)=\min\{n\geq0:a_n\geq y\}.
\]
For every fixed $J$, there are $K_J>0$ and $Y_J$ such that for
$y\geq Y_J$, with $x=x_J(y)$,
\begin{equation}\label{eq:ceil}
 \ceil{x-K_J\frac{x^{-(J+1)/2}}{\log x}}
 \ \leq I(y)\leq\
 \ceil{x+K_J\frac{x^{-(J+1)/2}}{\log x}}.
\end{equation}
Eventually the two displayed integers are equal or adjacent.
\end{theorem}
\begin{proof}
Appending a zero row and column injects the $n\times n$ class into the
$(n+1)\times(n+1)$ class. The matrix with a single one in the new
bottom-right position is outside the image, so $a_n$ is strictly
increasing. It is unbounded, for example because it contains every
permutation matrix. Thus $I(y)$ exists.

Put $\alpha=(J+1)/2$, $e(t)=t^{-\alpha}/\log t$, and
$z_n=x_J(a_n)$ for sufficiently large $n$. Equation
\eqref{eq:sequence-inverse} says $|z_n-n|\leq C e(n)$.
For large $y$ all relevant indices are in this range, and strict
monotonicity gives
\[
 I(y)=\min\{n:z_n\geq x\}.
\]
Since $e(n)\to0$, neighboring integer bounds first give
$I(y)=x+O(1)$. Explicitly, for all sufficiently large indices
$|z_n-n|<1$; hence indices below $x-1$ cannot cross, whereas any
integer at least $x+1$ does. Uniformly for integers $m=x+O(1)$,
$e(m)/e(x)\to1$, so $e(m)\leq2e(x)$ eventually.

Let $k=I(y)$. Then $x\leq z_k\leq k+Ce(k)$, so
$k\geq x-2Ce(x)$. Since $k$ is an integer,
$k\geq\ceil{x-K_Je(x)}$ whenever $K_J\geq2C$.
For the upper bound, choose $K_J>2C$ and put
$m=\ceil{x+K_Je(x)}$. Here $m=x+O(1)$, and
\[
 z_m\geq m-Ce(m)\geq x+(K_J-2C)e(x)\geq x.
\]
Thus $k\leq m$. This proves both sides, including at exact thresholds
$y=a_n$ and arbitrarily close to them. Finally $2K_Je(x)<1$
eventually, so the two ceilings differ by at most one.
\end{proof}

The enclosure is an asymptotic theorem with existential constants.
It is not a finite-input certified rounding algorithm. In general one
cannot round an asymptotic real inverse in a universally correct way:
$y$ can be arbitrarily close to a discrete threshold. Formula
\eqref{eq:inverse} alone therefore gives $I(y)$ equal to its displayed
real expression plus $O(1)$, while \eqref{eq:ceil} supplies the stronger
rounding-safe statement when the comparison and error allowance are
retained.

\section{Exact fixtures and reproducibility}\label{sec:checks}
The seventeen values in the inspected OEIS record, beginning at $n=0$,
are reproduced below as finite numeric fixtures. The companion checks
every value against the exact recurrence.
\begin{center}\small
\begin{tabular}{r r}
\toprule $n$ & $a_n$\\\midrule
0&1\\1&2\\2&16\\3&265\\4&7343\\5&304186\\6&17525812\\
7&1336221251\\8&129980132305\\9&15686404067098\\
10&2297230134084416\\11&400977650310256537\\
12&82188611938415464231\\13&19536244019455339261970\\
14&5328019975275896220786388\\
15&1651867356348327784988233291\\
16&577522171260292028520919811777\\\bottomrule
\end{tabular}
\end{center}
The following independent counting construction does not use the ODE.
For a fixed number $n$ of columns, let a state $(k,\ell)$ record the
numbers of columns currently of degree zero and one. A new row has
zero, one, or two ones. Its transitions, with multiplicities, are
\[
\begin{array}{c|c}
\text{new state}&\text{multiplicity}\\\hline
(k,\ell)&1\\
(k-1,\ell+1)&k\\
(k,\ell-1)&\ell\\
(k-2,\ell+2)&\binom{k}{2}\\
(k-1,\ell)&k\ell\\
(k,\ell-2)&\binom{\ell}{2}
\end{array}
\]
Start at $(n,0)$, repeat for $n$ rows, and sum the state weights.
Invalid states are omitted. A more literal check retains the entire
labeled column-degree tuple. Small binary matrices can also be
exhaustively checked by their $n^2$-bit masks. These constructions
supply checks with genuinely different state representations.

The source archive contains the LaTeX source, deterministic PDF and
ZIP builders, a manifest, finite numeric fixtures, and bounded
companion programs. The core exact arithmetic uses standard-library
Python integers and rational numbers; the finite symbolic
coefficient calculations also use only the standard library. The companion's usage guide specifies implemented input bounds,
commands, and precise test coverage. The code uses explicit exceptions
and test assertions that remain active with Python optimization.
No network access is needed to reproduce the packaged computations.

Finite exact agreement verifies implementations and identities, not
the asymptotic theorem. Conversely, the proof gives no effective
finite-$n$ error constant. Any decimal diagnostic in the companion is
uncertified floating-point evidence and must not be read as an interval
bound or proof. The release README records build requirements and
replay commands; the verification receipt reports the actual completed
checks.

\section{Questions for further research}\label{sec:future}
The exact bivariate function \eqref{eq:egf} suggests several concrete
extensions that are not proved here.

\paragraph{Rectangular aspect ratios.}
For $m\neq n$, the coefficient $m!n![x^my^n]F(x,y)$ invites a bivariate saddle analysis, or a uniform analysis of
Bessel functions whose order grows with the dimensions. A useful first target is a uniform leading
formula when $m/n$ stays in a compact subinterval of $(0,\infty)$,
followed by the transition as the aspect ratio approaches one on a
shrinking scale. The balanced Bessel diagonal alone does not supply
that uniformity.

\paragraph{Marked components and distributions.}
Separate variables can mark isolated vertices, balanced and unbalanced
paths, or cycles in the connected-component exponent. Differentiating
an appropriate marked coefficient expansion could yield means,
covariances, and limit laws for a uniformly chosen matrix. Such claims
require complex-uniform marked estimates and control of the
normalization; the unmarked expansion does not by itself prove them.

\paragraph{Effective errors and onset.}
Explicit constants in the Bessel remainder, arc estimates, and Taylor
majorant could turn the eventual two-ceiling theorem into a certified
finite-input enclosure. This would require a quantitative onset and
verified error allowances, rather than fitting the constants from
finite exact values.

\paragraph{Exponentially small terms.}
The present minor-arc bound is sufficient for every fixed algebraic
order, but does not classify secondary saddles or Stokes effects of
the Bessel factor. Determining late coefficient growth, optimal
truncation, and exponentially small contributions is a separate
problem. No complete transseries is supplied by the current proof.

\section{Literature coverage and claim boundaries}\label{sec:literature}
The inspected A197458 record defines the same square binary class and
provides a multivariate recurrence and Java implementation
\cite{oeis}. Its revision is number 25, dated 8 April 2020; the
source retrieval on 3 October 2026 found no asymptotic formula in that
record. Mathar's seven-page note gives recurrences, numerical tables,
and a Maple implementation; its Table~2 contains the relevant
rectangular family \cite{mathar}. The linked original discussion gives
the recurrence rather than an asymptotic \cite{mse}.

Closely related component EGFs are older. OEIS A062154 records a
bivariate EGF for matrices over $\{0,1,2\}$ with line sums one or two,
and it and the square companion A062156 cite Goulden and Jackson,
\emph{Combinatorial Enumeration}, Problem~3.4.15 \cite{companion}.
The printed exercise and its answer were not directly inspected.
Their exact scope remains an explicit source gap; the indirect
citation must not be converted into a claim about an unread theorem.
The exactly-two binary family A001499 is also classical, and the
retrieved OEIS entry records a leading equivalent and older
references \cite{regular}. This nearby class is not A197458, but it
precludes any suggestion that degree-two bipartite enumeration itself
is new.

Greenhill, McKay, and Wang's uniform asymptotics for prescribed sparse
row and column degree sequences include bounded degrees at most two
\cite{gmw}. Their theorem does not explicitly perform the sum over all
vectors defining $a_n$, but it provides classical machinery from
which a leading specialization can be approached. Caizergues and
de Panafieu give a general exact bounded-degree bipartite generating
function, allowing different degree weights on the two shores
\cite{cp}; this includes the finite degree set of the present class.
The retrieved paper does not state the scalar asymptotic proved here.
These facts warrant crediting general enumeration methods rather than
claiming a new enumeration principle.

The present work supplies an elementary explicit specialization,
a sharp all-fixed-orders calculation, four displayed corrections,
and an inverse with a safe discrete interpretation. A bounded search
of retrieved sources did not locate the particular scalar expansion;
that is not a historical-priority certificate. Older degree-two
terminology and the unread textbook exercise leave uncertainty. No
claim of first discovery is needed for any mathematical statement in
this report.

The primary analytic scope is also deliberately narrow: all orders
are fixed before taking the limit; no complete beyond-all-orders
transseries, convergent coefficient expansion, effective global
remainder constant, or universally correct rounded inverse is
asserted. The theorem and exact checks remain valid independently of
whether the scalar equivalent appears elsewhere in the literature.

\begin{thebibliography}{9}
\bibitem{oeis}
OEIS Foundation Inc., A197458, \emph{Number of $n\times n$ binary
matrices with at most two ones in each row and column}.
\url{https://oeis.org/A197458}.
Inspected record revision 25, 8 April 2020; retrieved 3 October 2026.
Official data record:
\url{https://github.com/oeis/oeisdata/blob/main/seq/A197/A197458.seq}.

\bibitem{mathar}
R. J. Mathar, \emph{The number of binary $n\times m$ matrices with at
most $k$ ones in each row or column}, 23 November 2014, seven pages.
\url{https://oeis.org/A247158/a247158.pdf}.

\bibitem{mse}
\emph{Filling out an $n\times n$ square grid with zeros and ones},
Math Stack Exchange question 72600 (2011), answer by joriki (2012).
\url{https://math.stackexchange.com/questions/72600}.
Original exposition of the recurrence; not a refereed publication.

\bibitem{companion}
OEIS Foundation Inc., A062154 and A062156.
\url{https://oeis.org/A062154}; \url{https://oeis.org/A062156}.
These entries cite I. P. Goulden and D. M. Jackson,
\emph{Combinatorial Enumeration}, Wiley, 1983, Problem 3.4.15.
The book exercise was not directly inspected.

\bibitem{regular}
OEIS Foundation Inc., A001499, the exactly-two line-sum family.
\url{https://oeis.org/A001499}.
Historical attributions here are via the entry; the cited original
Knuth solution was not directly inspected.

\bibitem{gmw}
C. Greenhill, B. D. McKay, and X. Wang,
\emph{Asymptotic enumeration of sparse 0--1 matrices with irregular
row and column sums}, 2006, Theorem 1.3.
Author-hosted text: \url{https://cs.anu.edu.au/~bdm/papers/irregular.pdf}.

\bibitem{cp}
E. Caizergues and \'E. de Panafieu,
\emph{Exact enumeration of graphs and bipartite graphs with degree
constraints}, EuroComb 2023, pp.~254--262, Theorem 2.
\url{https://doi.org/10.5817/CZ.MUNI.EUROCOMB23-035}.
Publisher text: \url{https://journals.muni.cz/eurocomb/article/view/35569}.
\end{thebibliography}
\end{document}
